
%Coreference resolution is one of the fundamental tasks in natural language processing. In this paper, we demonstrate that the common corpus for coreference resolution exhibits gender biases. As a result, the models trained on the corpus inherit these biases and make systematic errors on gender pronouns. 
%In this paper, we build a new winograd dataset which can be used as a benchmark to test if the coreference model exists gender bias issue.
%To mitigate the biases, we propose to augment the corpus with gender-reversed corpus such that the impact of the gender bias can be neutralized. We also propose to use the debiased word embeddings and the combination of these two methods to mitigate the bias in coreference resolution system. Experimental results show that the combination method shows the best calibration effect.
% To mitigate the biases, we proposed to augment the original corpus with a gender-reversed corpus, such that the impact of the gender bias can be neutralized. Experimental results show that the models trained on the augmented corpus achieve similar performance both on the original test corpus as well as the gender-reversed corpus. 

%can be adopted in many areas such as information extraction systems, question answering or machine translation. The broad  application  can be problematic if the underlying model presents  inappropriate societal biases.
%Taking gender bias as an example,  we demonstrate the OntoNote dataset is biased from different perspectives such as the gender ratio for the clusters or for different job titles. By building up the adversarial development dataset, we find the coreference resolution systems also exhibit the problem.  
% To help reduce the gender bias in the system, we retrain the model on a new adversarial dataset which can be obtained from the original training dataset. The result shows our ART algorithm can help to reduce the bias without affecting on the performance of the original models. 
Bias in NLP systems has the potential to not only mimic but also amplify stereotypes in society. 
For a prototypical problem, coreference, we provide a method for detecting such bias and show that three systems are significantly gender biased.
We also provide evidence that systems, given sufficient cues, can ignore their bias.
Finally, we present general purpose methods for making co-reference models more robust to spurious, gender-biased cues while not incurring significant penalties on their performance on benchmark datasets. 

%In this work, the corpus we use and the dataset we build is in English. Also here we consider the binary gender case. In the future it will be of much interest to work on some other languages such as Turkish where we cannot tell the gender from the singular third person pronouns. 

%We show that in three prototypical co-reference systems, gender bias from both supporting resources and training data significantly impact system performance, and 

%\my{impact statement here, because well, we have no space anywhere else...}
%To the best of our knowledge, our work is the first to demonstrate gender bias issues in coreference resolution and  propose a new dataset to verify this bias.
%The proposed algorithms are simple but effective.
%We are interested in extending this work by applying our methods to other tasks.%artificial intelligence tasks and designing new algorithms to mitigate the impact of bias in data by adding regularization in the training objective function.
% propose an approach to reduce this bias by generating a gender-reversed corpus. The proposed ART algorithm is simple but effective and does not require additional data annotations. We are interested in extending this work by applying ART to other artificial intelligence tasks and designing new algorithms to mitigate the impact of bias in data by adding regularization in the training objective function. 

\section*{Acknowledgement}
This work was supported in part by National Science Foundation Grant IIS-1760523, two NVIDIA GPU Grants, and a Google Faculty Research Award. We would like to thank Luke Zettlemoyer, Eunsol Choi, and Mohit Iyyer for helpful discussion and feedback.
